%(BEGIN_QUESTION)
% Copyright 2006, Tony R. Kuphaldt, released under the Creative Commons Attribution License (v 1.0)
% This means you may do almost anything with this work of mine, so long as you give me proper credit

En lastebil som veier 39,000 newton (N) kjører i 10 meter per sekund når clutchen ryker, slik at motoren kobles fra drivverket (girkasse, aksel, hjul osv.). Feilen skjer akkurat ved foten av en bratt bakke. Hvor mange meter (vertikalt) vil lastebilen rulle opp bakken uten motorkraft før den stopper, hvis vi ser bort fra all friksjon?

\vskip 10pt

Hvor høyt ville lastebilen rullet dersom den hadde kjørt dobbelt så fort?

\underbar{file i00431}
%(END_QUESTION)





%(BEGIN_ANSWER)

5.1 meter (målt vertikalt) ved starthastighet 10 m/s. Ved 20 m/s ville lastebilen ha oppnådd {\it fire ganger} så stor høydegevinst (20.4 meter).

\vskip 10pt

Siden potensiell energi ved stoppunktet i bakken skal være lik kinetisk energi i det clutchen ryker (forutsatt null energitap til friksjon), kan vi løse for $h$ enkelt:

$$E_p = mgh \hbox{\hskip 100pt} E_k = {1 \over 2}mv^2$$

$$mgh = {1 \over 2}mv^2$$

$$gh = {1 \over 2}v^2$$

$$h = {{1 \over 2}v^2 \over g}$$

$$h = {v^2 \over 2g}$$

Merk at bakkens helning ikke er oppgitt, fordi den er irrelevant for svaret på hvor stor vertikal høyde lastebilen får ved rulling. Hva {\it ville} bakkehelningen påvirke?

%(END_ANSWER)





%(BEGIN_NOTES)


%INDEX% Physics, energy, work, power: kinetic energy transformed to potential

%(END_NOTES)
